% HeatUp — npj Computational Materials
% Overleaf: upload main.tex + references.bib
% Compile: pdflatex → bibtex → pdflatex → pdflatex
%
% npj CompMat author guidelines (2024):
%   Article: ~4000 words main text, max 8 figures
%   Sections: Introduction | Results | Discussion | Methods
%   Abstract: ≤200 words, unstructured
%   References: numbered superscript (unsrtnat)
%   Figures: 85 mm (1 col) or 170 mm (2 col), 300 dpi min, sans-serif

\documentclass[twocolumn,10pt]{article}

%% ── Packages ─────────────────────────────────────────────────────────────────
\usepackage[utf8]{inputenc}
\usepackage[T1]{fontenc}
\usepackage{lmodern}
\usepackage[top=2.2cm,bottom=2.2cm,left=1.7cm,right=1.7cm,columnsep=5mm]{geometry}
\usepackage{amsmath,amssymb}
\usepackage{graphicx}
\usepackage{booktabs}
\usepackage[table]{xcolor}
\usepackage[hidelinks,colorlinks=true]{hyperref}
\usepackage{cleveref}
\usepackage{microtype}
\usepackage[font=small,labelfont=bf,justification=justified]{caption}
\usepackage{subcaption}
\usepackage{siunitx}
\usepackage{enumitem}
\usepackage[numbers,compress,super]{natbib}
\usepackage{listings}
\usepackage{authblk}
\usepackage{titlesec}
\usepackage{parskip}
\usepackage{fancyhdr}
\usepackage{abstract}

%% ── Colours ──────────────────────────────────────────────────────────────────
\definecolor{npjblue}{HTML}{1a4a8a}
\definecolor{codebg}{HTML}{F5F5F5}
\hypersetup{linkcolor=npjblue,citecolor=npjblue,urlcolor=npjblue}

%% ── Typography (Nature-style) ────────────────────────────────────────────────
\titleformat{\section}{\normalfont\bfseries\normalsize}{\thesection.}{0.5em}{}
\titleformat{\subsection}{\normalfont\bfseries\small}{\thesubsection.}{0.5em}{}
\setlength{\parskip}{3pt plus 1pt}
\setlength{\parindent}{0pt}

%% ── Abstract box ─────────────────────────────────────────────────────────────
\renewcommand{\abstractnamefont}{\bfseries\small}
\renewcommand{\abstracttextfont}{\small}
\setlength{\absleftindent}{0pt}
\setlength{\absrightindent}{0pt}

%% ── Code listings ────────────────────────────────────────────────────────────
\lstset{backgroundcolor=\color{codebg},basicstyle=\ttfamily\tiny,
  breaklines=true,frame=tb,framesep=2pt,rulecolor=\color{gray!30},
  language=Python,numbers=none,keywordstyle=\color{npjblue}\bfseries,
  commentstyle=\itshape\color{gray},xleftmargin=2pt,xrightmargin=2pt}

%% ── Math macros ──────────────────────────────────────────────────────────────
\newcommand{\kB}{k_{\mathrm{B}}}
\newcommand{\Fvib}{F_{\mathrm{vib}}}
\newcommand{\Fel}{F_{\mathrm{el}}}
\newcommand{\Fmag}{F_{\mathrm{mag}}}
\newcommand{\Fconf}{F_{\mathrm{conf}}}
\newcommand{\Ezero}{E_0}
\newcommand{\ehull}{\Delta E_{\mathrm{hull}}}
\newcommand{\TC}{T_{\mathrm{C}}}
\newcommand{\muB}{\mu_{\mathrm{B}}}
\newcommand{\vdos}{g(\omega)}
\newcommand{\meVat}{\,\mathrm{meV\,atom^{-1}}}
\newcommand{\tp}{\textsc{HeatUp}}

%% ── Page style ───────────────────────────────────────────────────────────────
\pagestyle{fancy}
\fancyhf{}
\fancyhead[L]{\small\textit{Lopez Alvarez \textbar{} HeatUp}}
\fancyhead[R]{\small\textit{npj Computational Materials}}
\fancyfoot[C]{\thepage}
\renewcommand{\headrulewidth}{0.4pt}

%% =============================================================================
\begin{document}
%% =============================================================================

%% ── Title block ──────────────────────────────────────────────────────────────
\title{\textbf{HeatUp}: a generalised Gibbs free-energy framework for\\
  temperature-aware phase stability in high-throughput materials discovery}

\author[1,*]{Cibran Lopez Alvarez}
\affil[1]{\small Department of Physics, Universitat Polit\`ecnica de Catalunya (UPC), Barcelona, Spain}
\affil[*]{\small\href{mailto:cibran.lopez@upc.edu}{cibran.lopez@upc.edu}}

\date{}
\maketitle
\thispagestyle{fancy}

%% ── Abstract ─────────────────────────────────────────────────────────────────
\begin{abstract}
Machine-learning interatomic potentials (MLIPs) have enabled high-throughput
computational screening of functional materials at near-first-principles accuracy,
producing large databases of candidates with promising predicted properties.
However, a predicted property alone does not guarantee that a material is
mechanically stable, vibrationally coherent, or thermodynamically accessible
at its operating temperature.
We introduce \tp{}, an open-source Python framework that closes this gap through a
sequential three-gate stability pipeline --- mechanical, vibrational, and thermodynamic
--- applied in order of increasing cost.
The thermodynamic gate constructs a fully temperature-dependent convex hull from a
generalised Gibbs free energy assembler combining anharmonic vibrational contributions
extracted from AIMD trajectories via the velocity autocorrelation function, electronic
free energies from the electronic density of states, magnetic free energies via the
Inden--Hillert model, and configurational entropy from site occupancies with a
Cowley short-range-order correction.
Critically, missing competing phases are generated automatically by exhaustive random
crystal-structure prediction across all 230 space groups using chemically motivated
stoichiometry hints to supplement the default search.
We demonstrate the framework on synthetic model systems that reproduce the physically
relevant behaviour of each contribution, and provide three self-contained Jupyter
notebooks for complete reproducibility.
\tp{} is designed as a drop-in stability filter for any MLIP-based discovery pipeline,
with a composable architecture that accommodates new contributions as MLIP capabilities evolve.
\end{abstract}

%% =============================================================================
\section*{Introduction}
%% =============================================================================

Universal MLIPs --- MACE-MP\cite{batatia2024maceuniversal}, CHGNet\cite{deng2023chgnet},
M3GNet\cite{chen2022universal}, SevenNet\cite{park2024sevennet} --- trained on millions
of DFT calculations now deliver near-first-principles accuracy for energies, forces,
and stresses at a cost several orders of magnitude lower.
Coupled with generative structure prediction\cite{xie2022crystal,yang2024mattersim}
and active-learning frameworks\cite{vandermause2020fly,lopezalvarez2024sse}, this has
produced large databases of candidates with predicted functional properties: high ionic
conductivity, target band gaps, optimal magnetic moments.

The central, underappreciated challenge is that a predicted property does not imply
that the hosting phase \emph{exists} under realistic conditions.
Three independent failure modes threaten any candidate:
mechanical instability (negative stiffness eigenvalue signals spontaneous
collapse\cite{born1954dynamical});
vibrational instability (anharmonic soft modes at operating temperature drive a
phase transition\cite{hellman2011lattice});
and thermodynamic instability (the material lies above the convex hull of competing
phases and will decompose\cite{sun2016thermodynamic}).

Standard databases --- the Materials Project\cite{jain2013commentary},
AFLOW\cite{curtarolo2012aflow} --- provide $T = 0\,\mathrm{K}$ hull data.
This approximation fails for functional materials operated at 700--1200\,K:
vibrational free-energy differences between phases reach $0.1$--$0.5\,\mathrm{eV\,atom^{-1}}$
at these temperatures\cite{wang2015design}, comparable to or exceeding typical hull
distances. Database coverage is also incomplete: novel compositions may have competing
phases not yet computed or generated\cite{sun2016thermodynamic}.

Here we introduce \tp{}, an open-source Python framework that systematically addresses
all three issues.
Every piece of data required for stability assessment is already produced as a byproduct
of MLIP validation workflows; \tp{} assembles it in the correct order, adds automated
phase generation, and returns a stability verdict at the operating temperature.

%% =============================================================================
\section*{Results}
%% =============================================================================

\subsection*{Sequential stability pipeline}

\tp{} evaluates three gates in order of increasing cost (Fig.~\ref{fig:pipeline}).
Failure at Gate~1 or Gate~2 terminates the pipeline; a warning propagates so that
borderline materials are fully characterised.

\textbf{Gate 1 --- Mechanical.}
The Born--Huang criterion\cite{born1954dynamical,mouhat2014necessary} requires all
six eigenvalues of the $6\times6$ Voigt stiffness tensor $C$ to be positive.
$C$ is computed from stress--strain finite differences using the MLIP.
Voigt-averaged bulk ($B$) and shear ($G$) moduli are also checked against configurable
thresholds. The calculation requires fewer than 50 single-point evaluations.

\textbf{Gate 2 --- Vibrational.}
Harmonic phonon calculations at 0\,K cannot capture anharmonic soft modes that develop
at finite temperature\cite{hellman2011lattice,muy2021phonon}.
Gate~2 uses the anharmonic VDOS extracted from the AIMD trajectory via the VACF:
\begin{equation}
  \vdos \propto \left|\,\mathcal{F}\!\Bigl[
    \textstyle\sum_{i,\alpha}\langle v_{i\alpha}(t)v_{i\alpha}(t+\tau)\rangle
  \Bigr]\,\right|. \label{eq:vacf}
\end{equation}
A material fails when the fraction $\zeta$ of integrated VDOS weight within
$|\omega| < 1\,\mathrm{meV}$ exceeds a configurable threshold (default: 8\%).
A minimum-frames check (default: 500 frames $\approx 10\,\mathrm{ps}$) guards
against spectral artefacts.
The VDOS is extracted from a trajectory already available from conductivity validation,
adding negligible overhead.

\textbf{Gate 3 --- Thermodynamic.}
The thermodynamic gate builds a temperature-dependent convex hull and evaluates the
hull distance $\ehull(T_{\rm op})$ at the specified operating temperature.

\subsection*{Generalised Gibbs free energy}

The total Gibbs free energy per atom is:
\begin{equation}
  G(T,P) = \Ezero + \Fvib(T) + \Fel(T) + \Fmag(T) + \Fconf(T) + PV,
  \label{eq:gibbs}
\end{equation}
assembled by the \texttt{GibbsAssembler} class, which accumulates contributions
additively and returns zero for absent data files.

\textbf{Anharmonic $\Fvib(T)$} uses the VDOS from Eq.~\eqref{eq:vacf}:
\begin{equation}
  \Fvib(T) = \int \vdos \bigl[{\tfrac{1}{2}}\hbar\omega
    + \kB T \ln(1 - e^{-\hbar\omega/\kB T})\bigr]\,\mathrm{d}\omega.
  \label{eq:fvib}
\end{equation}
A harmonic phonon DOS (from \texttt{phonons/dos.json}) is used as fallback when
no AIMD data are available.

\textbf{Electronic $\Fel(T)$} is the Mermin electronic free energy:
\begin{equation}
  \Fel(T) = U_{\rm el}(T) - T S_{\rm el}(T), \label{eq:fel}
\end{equation}
where $S_{\rm el} = {-\kB}\int g_{\rm el}[f\ln f+(1-f)\ln(1-f)]\,\mathrm{d}\varepsilon$
and the Fermi--Dirac chemical potential $\mu(T)$ is solved self-consistently to
conserve the electron number.
The electronic DOS $g_{\rm el}(\varepsilon)$ is read from
\texttt{electronic/edos.json}, which must be supplied from a DFT or tight-binding
calculation on the MLIP-relaxed structure.
\emph{Important limitation:} no deployed MLIP currently provides the electronic
DOS from first principles alone. This is an active research frontier; the
\texttt{GibbsAssembler} is designed to accommodate electron-density aware
MLIPs\cite{batatia2022mace} without interface changes once they become available.
For insulators with band gap $>1\,\mathrm{eV}$, $\Fel(T)$ is negligible
at $T < 2000\,\mathrm{K}$; for metals it contributes $5$--$50\meVat$.

\textbf{Magnetic $\Fmag(T)$} uses the Inden--Hillert
model\cite{inden1976role,dinsdale1991sgte} with exchange parameters from
\texttt{magnetic/moments.json}:
\begin{equation}
  \Fmag(T) = -\kB\TC\,f(\tau)\,\ln(\muB+1), \label{eq:fmag}
\end{equation}
where $\tau=T/\TC$ and $f(\tau)$ is the Inden polynomial.
This is quantitatively accurate for 3d transition metals ($\pm 10\meVat$),
but gives errors of $20$--$50\meVat$ for rare-earth $f$-electron
compounds\cite{barreda2022anharmonic}, where a full Heisenberg Hamiltonian
parameterised from MLIP spin dynamics would be needed --- not yet available in
deployed potentials.

\textbf{Configurational $\Fconf(T)$} from site occupancies in
\texttt{disorder/site\_occupancies.json}:
\begin{equation}
  S_{\rm conf} = -\frac{\kB}{N}\sum_i\sum_j x_{ij}\ln x_{ij}. \label{eq:sconf}
\end{equation}
A Warren--Cowley short-range-order (SRO) correction is applied when
\texttt{sro\_parameters.json} is present:
$S_{\rm conf}^{\rm corr} = \alpha_{\rm SRO}\,S_{\rm conf}^{\rm ideal}$, where
$\alpha_{\rm SRO}$ is derived from the AIMD radial distribution function.
This correction is critical for superionic conductors where SRO reduces the
ideal entropy by 30--50\%\cite{barreda2022anharmonic}.

\textbf{Pressure--volume $PV$} uses a Birch--Murnaghan
EOS\cite{birch1947finite} from \texttt{equation\_of\_state/eos.json}.
A full quasi-harmonic treatment is planned but not yet implemented.

\subsection*{Automated competing-phase generation}

For a target material with element set $\{E_1,\ldots,E_n\}$, \tp{} collects all
competing phases from the database and candidate tree, then generates missing
polymorphs via PyXtal\cite{fredericks2021pyxtal} across all 230 ITA space groups.
A \texttt{stoichiometry\_hints.json} file adds non-trivial stoichiometries for
common sub-systems: for example, the Li--P--S sub-system automatically includes
Li$_3$PS$_4$, Li$_2$PS$_3$, Li$_4$PS$_4$, and Li$_7$PS$_6$ in addition to the
1:1:1 default, covering NASICON, argyrodite, and LGPS-type systems.
Competing phases that already have AIMD trajectories are automatically upgraded
to anharmonic free energies; a per-phase log records which source was used.

\subsection*{Model demonstrations}

Figure~\ref{fig:gibbs} shows the free-energy decomposition for a model
lithium-halide system.
The anharmonic correction to $\Fvib$ reaches $-15\meVat$ at 1200\,K;
the configurational term (partial Li-site disorder) contributes $-100\meVat$.
The electronic term (insulating gap) is negligible, consistent with the
expectation that $\Fel$ matters only for metallic competing phases.

Figure~\ref{fig:hull} shows the temperature evolution of the convex hull for a
model binary Li--X system.
Two phases marked metastable at $T = 0$ reach the hull at 1200\,K due to larger
vibrational entropy and partial disorder.
A static 0\,K screen would incorrectly classify these as unstable.

Figure~\ref{fig:vdos} illustrates Gate~2 detection.
In panel (c), $\zeta = 12.4\%$ exceeds the default 8\% threshold, correctly
triggering pipeline termination.

Figure~\ref{fig:coverage} quantifies the impact of PyXtal generation on hull
distances for 80 synthetic candidates.
Including generated phases increases the median hull distance by $\sim30\meVat$
and reclassifies 15\% of materials as less stable, demonstrating that incomplete
databases systematically overestimate thermodynamic stability.

%% =============================================================================
\section*{Discussion}
%% =============================================================================

\subsection*{The stability gap in the MLIP era}

Universal MLIPs have removed the $E_0$ computation bottleneck from high-throughput
screening.
The new bottleneck is stability validation at realistic operating conditions.
\tp{} addresses this directly: every contribution to $G(T,P)$ in
Eq.~\eqref{eq:gibbs} is assembled from data already generated by the MLIP
validation workflow, at negligible marginal cost.
The sequential gating design ensures that the most expensive analysis --- hull
construction and phase generation --- is only performed for materials that survive
the cheaper mechanical and vibrational screens.

\subsection*{What is and is not currently achievable}

Table~\ref{tab:contributions} summarises the current feasibility of each
contribution. $\Fvib(T)$, $\Fconf(T)$, and $PV$ are fully computable from MLIP
data alone. $\Fmag(T)$ is accessible for 3d systems from DFT magnetic moments
but requires a separate DFT calculation; the Inden--Hillert model is insufficient
for $f$-electron compounds. $\Fel(T)$ requires an electronic DOS not available
from any deployed MLIP: obtaining it currently requires a DFT single-point
on the MLIP-relaxed geometry.

This situation is well-defined and temporary.
Electron-density and spin-density aware equivariant
models\cite{batatia2022mace} are an active research priority, and \tp{}'s
\texttt{GibbsAssembler} architecture accommodates them without interface changes.
Any function with signature \texttt{fn(sym\_dir, T\_array)\,$\to$\,np.ndarray}
can be registered and will be called automatically with the appropriate data.

\subsection*{Generality beyond solid-state electrolytes}

Although the motivating application is SSE discovery\cite{lopezalvarez2024sse},
the framework is material-class agnostic.
For thermoelectrics, a carrier-concentration-dependent $\Fel$ can be registered.
For high-entropy oxides, multi-sublattice $\Fconf$ with full SRO treatment.
For magnetic shape-memory alloys, a martensitic transformation free energy.
For high-pressure synthesis, a volume-grid QHA treatment of $PV$.
The four-line extension API --- register, implement, go --- is demonstrated in
Notebook~2 using a Schottky defect contribution as a minimal example.

\subsection*{Limitations and future directions}

\emph{Trajectory length.} Reliable VACF spectral resolution requires
$\geq 500$ production frames ($\approx 10\,\mathrm{ps}$); very low-frequency
modes need $\geq 100\,\mathrm{ps}$. \tp{} now emits a configurable warning
(\texttt{cfg.VIB\_MIN\_FRAMES}) when this is not met.

\emph{Short-range order.} The Warren--Cowley correction addresses the leading-order
SRO effect on $\Fconf$; multi-body and longer-range correlations require a
cluster-expansion treatment\cite{vandermause2020fly} beyond the current scope.

\emph{Quasi-harmonic pressure.} The current $PV$ uses a static BM EOS.
A volume-grid AIMD series would enable full QHA; automated support is planned.

\emph{Rare-earth magnetism.} The Inden--Hillert model is insufficient for $f$-electron
systems. A Heisenberg MC approach requires spin-adiabatic MLIPs not yet available.

\emph{Electronic DOS from MLIPs.} This is the highest-priority target for future
development, gating the use of $\Fel$ for metallic compounds without DFT.

%% =============================================================================
\section*{Methods}
%% =============================================================================

\subsection*{VACF and VDOS computation}

Velocities are estimated from PBC-aware finite differences of fractional coordinates
(minimum-image convention), converted to Cartesian \AA\,fs$^{-1}$.
The full two-sided VACF is computed via \texttt{numpy.correlate(...,'full')},
mean-subtracted before FFT to suppress the DC artefact.
The frequency axis uses $h/e = 4.1357\,\mathrm{meV\,fs}$.
VDOS results from multiple AIMD temperatures are averaged on a common grid.

\subsection*{Chemical potential solver}

$\mu(T)$ in $\Fel(T)$ is determined by bisection, conserving valence electrons
per atom to $< 10^{-6}$ at each temperature.

\subsection*{Inden--Hillert polynomial}

The dimensionless $f(\tau)$ uses CALPHAD normalisation constants
$D_{\rm below} = 0.6549$ and $D_{\rm above} = 0.6509$
from Ref.\,\citen{dinsdale1991sgte}.

\subsection*{Birch--Murnaghan EOS}

Third-order BM parameters are fitted by nonlinear least-squares
(\texttt{scipy.optimize.curve\_fit}).

\subsection*{Convex hull}

The pymatgen \texttt{PhaseDiagram} engine\cite{ong2013python} constructs the hull
at $T \in \{0,50,\ldots,1500\}\,\mathrm{K}$.

\subsection*{Thresholds}

All stability thresholds (Table~\ref{tab:thresholds}) are centralised in
\texttt{heatup/config.py} and overridable at runtime.

%% =============================================================================
%% Tables
%% =============================================================================

\begin{table}[ht]
\caption{\textbf{Free-energy contributions and current data requirements.}
  ``MLIP only'' indicates the contribution is computable without DFT.
  ``+ DFT'' requires a single-point DFT calculation.
  ``Not yet'' indicates a fundamental current limitation.}
\label{tab:contributions}
\small
\begin{tabular}{lllp{2.8cm}}
\toprule
Contribution & Data source & Status & Notes \\
\midrule
$\Fvib^{\rm anh}$ & AIMD VACF & MLIP only & Default \\
$\Fvib^{\rm harm}$ & Phonon DOS & MLIP only & Fallback \\
$\Fconf$ & Site occupancies & MLIP only & + SRO corr. \\
$PV$ & BM EOS & MLIP only & Static $T$ \\
$\Fel$ & Electronic DOS & + DFT & No MLIP yet \\
$\Fmag$ (3d) & DFT moments & + DFT & Inden--Hillert \\
$\Fmag$ ($f$) & Heisenberg MC & Not yet & Awaits spin MLIPs \\
$\Fvib$ QHA & Vol-grid AIMD & Planned & Future release \\
\bottomrule
\end{tabular}
\end{table}

\begin{table}[ht]
\caption{\textbf{Default stability thresholds in \tp{}}.
  All are overridable via \texttt{heatup/config.py}.}
\label{tab:thresholds}
\small
\begin{tabular}{lll}
\toprule
Parameter & Default & Unit \\
\midrule
\texttt{MECH\_BORN\_EIGENVALUE\_FAIL} & 0 & GPa \\
\texttt{MECH\_BULK\_WARN} & 10 & GPa \\
\texttt{MECH\_SHEAR\_WARN} & 5 & GPa \\
\texttt{VIB\_ZERO\_WINDOW} & 1 & meV \\
\texttt{VIB\_ZERO\_FRAC\_WARN} & 2 & \% \\
\texttt{VIB\_ZERO\_FRAC\_FAIL} & 8 & \% \\
\texttt{VIB\_MIN\_FRAMES} & 500 & frames \\
\texttt{THERMO\_HULL\_WARN} & 100 & meV/atom \\
\bottomrule
\end{tabular}
\end{table}

%% =============================================================================
%% Figures (replace \fbox+\parbox with \includegraphics in final version)
%% =============================================================================

\begin{figure*}[ht]
\centering
\fbox{\parbox{\textwidth}{\centering\vspace{1.2cm}
  \textbf{[Fig. 1: Pipeline schematic --- double column, 170 mm]}\\[0.3em]
  Generated by \texttt{plotting\_utils.fig1\_pipeline\_schematic()}.
  Sequential gates (Mechanical $\to$ Vibrational $\to$ Thermodynamic)
  with fail arrows, GibbsAssembler inputs, and PyXtal loop.
  \vspace{1.2cm}}}
\caption{\textbf{The \tp{} sequential stability pipeline.}
  Candidates enter from the left and traverse three gates in order of increasing
  cost.
  Gate~1 checks Born--Huang stability from the MLIP elastic tensor.
  Gate~2 inspects the anharmonic VDOS for soft-mode weight near $\omega=0$.
  Gate~3 builds the temperature-dependent convex hull from the generalised Gibbs
  free energy, with PyXtal automatically supplying missing competing phases.
  A failure terminates the pipeline; a warning propagates.
  All thresholds are configurable; see Table~\ref{tab:thresholds}.}
\label{fig:pipeline}
\end{figure*}

\begin{figure}[ht]
\centering
\fbox{\parbox{\columnwidth}{\centering\vspace{1.0cm}
  \textbf{[Fig. 2: Gibbs decomposition --- single column, 85 mm]}\\[0.2em]
  (a) Stacked contributions; (b) Harmonic--anharmonic difference.\\
  \texttt{fig2\_gibbs\_decomposition()}
  \vspace{1.0cm}}}
\caption{\textbf{Generalised Gibbs free-energy decomposition.}
  (a) Stacked contributions $\Fvib$, $\delta\Fvib^{\rm anh}$, $\Fel$, $\Fconf$
  to $G(T)-E_0$ for a model lithium-halide system.
  The configurational term and anharmonic vibrational correction dominate at
  $T_{\rm op} = 1200\,\mathrm{K}$ (dotted line).
  (b) Harmonic--anharmonic difference $\delta\Fvib$: the anharmonic correction
  reaches $-15\meVat$ at 1200\,K.}
\label{fig:gibbs}
\end{figure}

\begin{figure*}[ht]
\centering
\fbox{\parbox{\textwidth}{\centering\vspace{1.0cm}
  \textbf{[Fig. 3: Hull evolution --- double column, 3 panels]}\\[0.2em]
  $T = 0$, 600, 1200\,K; starred phases metastable at 0\,K, on hull at 1200\,K.\\
  \texttt{fig3\_hull\_evolution()}
  \vspace{1.0cm}}}
\caption{\textbf{Temperature evolution of the convex hull.}
  Formation free energies for a model binary Li--X system at 0, 600, and
  1200\,K.
  Phases Li$_3$X$_2^*$ and Li$_2$X$_3^\dagger$ are metastable at $T=0$
  but reach the hull at $T = 1200\,\mathrm{K}$ due to larger vibrational entropy
  and partial disorder.
  Green: on-hull. Orange: metastable. Stars: PyXtal-generated.
  A static 0\,K screen would incorrectly discard both starred phases.}
\label{fig:hull}
\end{figure*}

\begin{figure}[ht]
\centering
\fbox{\parbox{\columnwidth}{\centering\vspace{1.0cm}
  \textbf{[Fig. 4: VDOS --- single column, 3 panels]}\\[0.2em]
  (a) Harmonic DOS; (b) Anharmonic VDOS; (c) Soft-mode material.\\
  \texttt{fig4\_vdos\_comparison()}
  \vspace{1.0cm}}}
\caption{\textbf{Vibrational stability assessment (Gate~2).}
  (a) Harmonic phonon DOS.
  (b) Anharmonic VDOS from VACF: thermally broadened, red-shifted peaks.
  (c) Vibrationally unstable material: soft-mode fraction
  $\zeta = 12.4\%$ within $|\omega| < 1\,\mathrm{meV}$ (red band)
  exceeds the 8\% failure threshold, terminating the pipeline.}
\label{fig:vdos}
\end{figure}

\begin{figure}[ht]
\centering
\fbox{\parbox{\columnwidth}{\centering\vspace{1.0cm}
  \textbf{[Fig. 5: Coverage impact --- single column, 2 panels]}\\[0.2em]
  (a) Paired scatter; (b) Reclassification matrix.\\
  \texttt{fig5\_coverage\_impact()}
  \vspace{1.0cm}}}
\caption{\textbf{Effect of systematic phase generation on hull distances.}
  (a) Hull distance $\ehull$ before vs after adding PyXtal-generated phases
  for 80 model candidates.
  (b) Reclassification matrix: 15\% of materials are reclassified as
  ``metastable'' or ``unstable'' once the competing-phase space is systematically
  completed, demonstrating that incomplete databases overestimate stability.}
\label{fig:coverage}
\end{figure}

%% ── Data and code availability ───────────────────────────────────────────────
\section*{Data availability}
All data required to reproduce the figures are generated analytically within the
companion Jupyter notebooks included in the repository.
No proprietary data are used.

\section*{Code availability}
\tp{} is available at \url{https://github.com/cibranlopez/heatup} under
the MIT licence (Zenodo DOI: \texttt{10.5281/zenodo.XXXXXXX}).
The version corresponding to this paper is tagged \texttt{v0.1.0}.

\section*{Acknowledgements}
The author thanks colleagues at UPC for stimulating discussions on active learning
and thermodynamic stability, and the developers of
pymatgen\cite{ong2013python}, PyXtal\cite{fredericks2021pyxtal},
and MACE\cite{batatia2022mace} for excellent open-source tools.

\section*{Author contributions}
C.L.A.\ conceived the framework, developed the software, performed the analysis,
and wrote the manuscript.

\section*{Competing interests}
The author declares no competing interests.

\bibliographystyle{unsrtnat}
\bibliography{references}

\end{document}
